% Radar (polar) plot: scaffold comparison across seven normalised axes.
% Adapted from raw/cybench-scaffolds-FINAL.pdf Figure 2; colours rebased on the
% paper palette and CSI::GCAI Cmd Accuracy recomputed using the recovered
% command count from Table 1 (commands=2491, errors=340 -> 86.4%).
%
% Axis values are normalised to [0,1]. Best per axis = 1.0.
%   Solve Rate     = Solved / 33
%   Speed          = min(WallTime) / WallTime         (CAI is reference)
%   Cost Eff       = min(TotalCost) / TotalCost       (CAI is reference)
%   Cmd Accuracy   = 1 - Errors/Commands              (Codex is reference)
%   Token Eff      = min(TotalTokens) / TotalTokens   (CAI is reference)
%   Breadth        = #difficulty-tiers covered / 5    (all four cover all five)
%   Simplicity     = min(LOC) / LOC                   (GCAI 20.4k LOC reference)
%
% Per-scaffold values:
%        Claude  Codex   GCAI    CAI
% Solve  0.545   0.606   0.515   0.576
% Speed  0.111   0.146   0.107   1.000
% CostEf 0.082   0.159   0.115   1.000
% CmdAcc 0.631   0.946   0.864   0.692
% TokEf  0.016   0.035   0.013   1.000
% Breadt 1.000   1.000   1.000   1.000
% Simpl  0.007   0.043   1.000   0.070
\begin{tikzpicture}[scale=2.4, every node/.style={font=\small\sffamily}]
  % Seven axes at angles 90, 90-360/7, ...
  \def\Naxes{7}
  \pgfmathsetmacro{\angstep}{360/\Naxes}

  % Background grid: 4 concentric circles + axes
  \foreach \r in {0.25,0.5,0.75,1.0}
    \draw[gray!30, line width=0.3pt] (0,0) circle (\r);
  \foreach \i in {1,...,7}{
    \pgfmathsetmacro{\ang}{90 - (\i-1)*\angstep}
    \draw[gray!40, line width=0.3pt] (0,0) -- (\ang:1.05);
  }
  % Tick labels on the upper-right axis (i=1, top)
  \node[gray!70, font=\tiny\sffamily, anchor=south] at (90:0.27) {25\%};
  \node[gray!70, font=\tiny\sffamily, anchor=south] at (90:0.52) {50\%};
  \node[gray!70, font=\tiny\sffamily, anchor=south] at (90:0.77) {75\%};
  \node[gray!70, font=\tiny\sffamily, anchor=south] at (90:1.02) {100\%};

  % Axis labels
  \def\axlabels{
    1/Solve Rate/0.50,
    2/Speed/0.85,
    3/Cost Eff./0.85,
    4/Cmd Acc./0.85,
    5/Token Eff./0.85,
    6/Breadth/0.85,
    7/Simplicity/0.50}
  \foreach \i/\lab/\lr in \axlabels {
    \pgfmathsetmacro{\ang}{90 - (\i-1)*\angstep}
    \node[font=\footnotesize\sffamily, color=cai_dark] at (\ang:1.18) {\lab};
  }

  % Polygons per scaffold. Order of axes: Solve, Speed, CostEf, CmdAcc, TokEf, Breadth, Simplicity
  % Helper: coordinate at axis i, value v.
  % \pgfmathsetmacro{\ang}{90 - (i-1)*60} -> (\ang:v)
  % Claude
  \def\polyClaude{
    (90:0.545) -- (38.57:0.111) -- (-12.86:0.082) -- (-64.29:0.631) --
    (-115.71:0.016) -- (-167.14:1.0) -- (-218.57:0.007) -- cycle}
  \def\polyCodex{
    (90:0.606) -- (38.57:0.146) -- (-12.86:0.159) -- (-64.29:0.946) --
    (-115.71:0.035) -- (-167.14:1.0) -- (-218.57:0.043) -- cycle}
  \def\polyGCAI{
    (90:0.515) -- (38.57:0.107) -- (-12.86:0.115) -- (-64.29:0.864) --
    (-115.71:0.013) -- (-167.14:1.0) -- (-218.57:1.0) -- cycle}
  \def\polyCAI{
    (90:0.576) -- (38.57:1.0) -- (-12.86:1.0) -- (-64.29:0.692) --
    (-115.71:1.0) -- (-167.14:1.0) -- (-218.57:0.07) -- cycle}

  % Drawing order = bottom-to-top in z-stack. We draw the two cai_primary-derived
  % scaffolds first, then Codex, and Claude last so the vendor-coded red/blue
  % polygons sit on top and remain easy to identify.
  \fill[gcai_color, opacity=0.30]  \polyGCAI;
  \draw[gcai_color!75!black, line width=0.9pt] \polyGCAI;
  \fill[cai_orange, opacity=0.28]  \polyCAI;
  \draw[cai_orange!75!black, line width=0.9pt] \polyCAI;
  \fill[codex_color, opacity=0.28] \polyCodex;
  \draw[codex_color!75!black, line width=1.1pt] \polyCodex;
  \fill[cc_color, opacity=0.28]    \polyClaude;
  \draw[cc_color!75!black, line width=1.1pt] \polyClaude;

  % Vertex markers
  \foreach \i/\v/\c in {
    1/0.545/cc_color,2/0.111/cc_color,3/0.082/cc_color,4/0.631/cc_color,5/0.016/cc_color,6/1.0/cc_color,7/0.007/cc_color}{
    \pgfmathsetmacro{\ang}{90 - (\i-1)*\angstep}
    \fill[\c] (\ang:\v) circle (0.014);
  }
  \foreach \i/\v/\c in {
    1/0.606/codex_color,2/0.146/codex_color,3/0.159/codex_color,4/0.946/codex_color,5/0.035/codex_color,6/1.0/codex_color,7/0.043/codex_color}{
    \pgfmathsetmacro{\ang}{90 - (\i-1)*\angstep}
    \fill[\c] (\ang:\v) circle (0.014);
  }
  \foreach \i/\v/\c in {
    1/0.515/gcai_color,2/0.107/gcai_color,3/0.115/gcai_color,4/0.864/gcai_color,5/0.013/gcai_color,6/1.0/gcai_color,7/1.0/gcai_color}{
    \pgfmathsetmacro{\ang}{90 - (\i-1)*\angstep}
    \fill[\c] (\ang:\v) circle (0.014);
  }
  \foreach \i/\v/\c in {
    1/0.576/cai_orange,2/1.0/cai_orange,3/1.0/cai_orange,4/0.692/cai_orange,5/1.0/cai_orange,6/1.0/cai_orange,7/0.07/cai_orange}{
    \pgfmathsetmacro{\ang}{90 - (\i-1)*\angstep}
    \fill[\c] (\ang:\v) circle (0.014);
  }

  % Legend (bottom, single row)
  \begin{scope}[shift={(0,-1.7)}]
    \fill[cc_color, opacity=0.5]   (-2.0,0) rectangle (-1.85,0.1);
    \node[anchor=west, font=\footnotesize\sffamily, color=cc_color] at (-1.83,0.05) {CSI::Claude};
    \fill[codex_color, opacity=0.5] (-0.7,0) rectangle (-0.55,0.1);
    \node[anchor=west, font=\footnotesize\sffamily, color=codex_color] at (-0.53,0.05) {CSI::Codex};
    \fill[gcai_color, opacity=0.5] (0.55,0) rectangle (0.7,0.1);
    \node[anchor=west, font=\footnotesize\sffamily, color=gcai_color] at (0.72,0.05) {CSI::GCAI};
    \fill[cai_orange, opacity=0.5] (1.85,0) rectangle (2.0,0.1);
    \node[anchor=west, font=\footnotesize\sffamily, color=cai_orange] at (2.02,0.05) {CSI::CAI};
  \end{scope}
\end{tikzpicture}
